% Generated draft: requires mathematical and source review.
\documentclass[11pt]{article}
\usepackage[margin=25mm]{geometry}
\usepackage{fontspec}
\setmainfont{Latin Modern Roman}
\usepackage{amsmath,amssymb,mathtools,unicode-math}
\setmathfont{Latin Modern Math}
\usepackage{enumitem,needspace,xurl,hyperref}
\setlength{\parindent}{0pt}
\setlength{\parskip}{5pt}
\setlength{\emergencystretch}{3em}
\newcommand{\N}{\mathbb{N}}
\newcommand{\Z}{\mathbb{Z}}
\newcommand{\Q}{\mathbb{Q}}
\newcommand{\R}{\mathbb{R}}
\newcommand{\C}{\mathbb{C}}
\newcommand{\F}{\mathbb{F}}
\newcommand{\A}{\mathbb{A}}
\newcommand{\PP}{\mathbb{P}}
\newcommand{\E}{\mathbb{E}}
\newcommand{\eps}{\varepsilon}
\newcommand{\dd}{\,\mathrm{d}}
\begin{document}
\section*{UnsolvedMath ID: 30003308}
\textbf{Source problem code:} catalogue entry
\subsection{Definitions and mathematical statement}
Let \(a_1,\ldots,a_5\) be positive integers. Interpret the source's phrase ``relatively prime'' collectively as \(\gcd(a_1,\ldots,a_5)=1\); pairwise coprimality is not explicitly required. Put
\[
S=\left\{\sum_{i=1}^{5}n_i a_i:n_i\in\Z_{\geq 0}\right\}.
\]
For a field \(k\), the associated affine monomial curve is
\[
C_a=\operatorname{Spec}R_a\subseteq\mathbb A_k^5,\qquad
R_a=k[t^{a_1},\ldots,t^{a_5}]
\cong k[x_1,\ldots,x_5]/\ker\!\left(x_i\mapsto t^{a_i}\right).
\]
Here Gorenstein means that every local ring of \(R_a\) is Gorenstein, where a Noetherian local ring is Gorenstein if it has finite injective dimension as a module over itself. The source does not specify the ground field, its characteristic, or whether Gorensteinness is intended globally or at the origin; this gives the standard coordinate-ring interpretation, subject to review.

A conventional meaning of a principal matrix is the following. For each \(i\), let
\[
c_i=\min\left\{c\in\Z_{>0}:ca_i\in\sum_{j\ne i}\Z_{\geq 0}a_j\right\}.
\]
Choose coefficients \(m_{ij}\in\Z_{\geq 0}\) satisfying
\[
c_i a_i=\sum_{j\ne i}m_{ij}a_j.
\]
The matrix \(M=(m_{ij})\in\Z^{5\times5}\), with diagonal entries \(m_{ii}=-c_i\), then satisfies
\[
M(a_1,\ldots,a_5)^{\mathsf T}=0.
\]
The off-diagonal coefficients can be nonunique. The supplied record does not define ``principal matrix,'' so this conventional definition and its sign convention require confirmation against the report.

\textbf{Requested characterization.} For all five-tuples satisfying the intended hypotheses, find an explicit necessary and sufficient criterion, expressed in terms of their principal matrix, for \(C_a\) to be Gorenstein. A precise criterion must specify how it treats different principal matrices of the same tuple. The record does not require the five integers to be distinct or to form a minimal generating set of \(S\); neither restriction is imposed here.

The supplied dated assessment says that a principal-matrix criterion is available for decomposable length-five Gorenstein sequences, while the requested characterization for all length-five sequences is not established. The record supplies neither the definition of ``decomposable'' nor the restricted criterion. This qualification is retained as supplied curation, not as an independently verified mathematical or literature claim.
\subsection{Short English statement}
Find a necessary and sufficient condition on the principal matrix of five relatively prime positive integers for their affine monomial curve to be Gorenstein. The source leaves the ground field, minimality of the generators, and treatment of nonunique principal matrices unspecified. Its supplied assessment reports a partial result for decomposable sequences; the general characterization remains the requested target.
\subsection{Sources}
\begin{itemize}
\item Supplied local UnsolvedMath registry record, permanent ID \(30003308\), title ``Gorenstein Monomial Curves from Principal Matrices,'' source problem code \texttt{catalogue entry}, in the collection ``Oberwolfach Reports Open Problems.'' No individual proposer is supplied; the recorded proposed year is 2016.
\item Supplied report citation: \emph{Asymptotic Phenomena in Local Algebra and Singularity Theory} (2017), Oberwolfach Reports, pp.~3199--3259, DOI \texttt{10.4171/owr/2016/57}. \url{https://doi.org/10.4171/owr/2016/57}. The record attributes the problem to workshop open-problem discussions documented by this report.
\item Supplied literature-source entry: ``OWR 57 (2017),'' \url{https://doi.org/10.4171/owr/2016/57}. Literature checked at: \texttt{2026-08-21}. This exact date records supplied dated curation, not independent verification. The supplied status is \texttt{partially\_solved}; its assessment reports a criterion for decomposable length-five Gorenstein sequences and says that the requested criterion for all length-five sequences is not established. The background explicitly requests independent verification of time-sensitive status and cited claims.
\item The supplied statement status is \texttt{corrected\_verified}, with a note that only notation, typography, and TeX delimiters were cleaned up. This is registry metadata, not independent mathematical verification of this draft.
\end{itemize}
\end{document}
